\section{The free part: \texorpdfstring{$\HF$}{HF} and the auditable axioms}\label{sec:free}

We now run the instrument from nothing and read off what it supplies. Three results carry the
section: the packaging map is an idempotent whose objects are the hereditarily finite sets
(\S\ref{ssec:p5}); closing the empty family under the finite packaging operations gives exactly
$\HF$, which satisfies every axiom of $\mathsf{ZFC}$ except Infinity (\S\ref{ssec:hf}); and on
finite graphs Foundation is the same as the existence of a rank audit (\S\ref{ssec:foundation}).
Table~\ref{tab:freepart} summarizes the outcome.

\subsection{The set universe as a packaging quotient}\label{ssec:p5}

\begin{theorem}[Mostowski collapse is packaging; \inhouse]\label{thm:p5}
The packaging map $E\colon Z\to Z$ of Definition~\ref{def:E} is idempotent, and
\[
  \Fix(E)\;=\;\im(E)\;=\;\{C(a): a\in\HF\}.
\]
The restricted lens $f\restr\Fix(E)\colon \Fix(E)\to\HF$ is a bijection with inverse $a\mapsto
C(a)$. With $i\colon\Fix(E)\hookrightarrow Z$ the inclusion and $r\colon Z\to\Fix(E)$,
$r(P):=E(P)$,
\[
  r\circ i=\mathrm{id}_{\Fix(E)},\qquad i\circ r=E,
\]
so $\Fix(E)$ is a system of representatives for the lens quotient $Z/{\sim}$.
\end{theorem}

\begin{proof}
Since $f(C(a))=a$ (Lemma~\ref{lem:L2}), $E(E(P))=C(f(C(f(P))))=C(f(P))=E(P)$. Idempotence gives
$\im(E)\subseteq\Fix(E)$, the reverse inclusion is trivial, and $\im(E)=\{C(a):a\in\HF\}$ by
definition. The restricted lens is surjective because $f(C(a))=a$ and injective because $P=E(P)=
C(f(P))$ on $\Fix(E)$. The splitting equations are immediate, and $[P]\mapsto E(P)$ is a
bijection from $Z/{\sim}$ onto $\Fix(E)$ because $E(P)\sim P$. This is the framework's
``idempotents split'' lemma~\citep{SBTFoundationsI} for the membership instrument; the
underlying collapse of a well-founded extensional relation onto a transitive set is the Mostowski
lemma~\citep{Mostowski1949}.
\end{proof}

\begin{remark}[The construction already in use]\label{rem:mathlib}
Theorem~\ref{thm:p5} describes how the universe of sets is commonly built in proof assistants.
The \texttt{PSet}/\texttt{ZFSet} development of \citet{Mathlib2020} defines a pre-set as a
type-indexed family of pre-sets---a well-founded membership tree with arbitrary branching---and a
set as a class of pre-sets under extensional equivalence. That is the quotient of presentations
by the lens relation $\sim$, and Theorem~\ref{thm:p5} is its finite, well-founded instance. The
Lean development accompanying this paper builds the finite version from scratch, as a quotient of
finite membership trees, without choosing representatives (\S\ref{sec:mech}).
\end{remark}

\begin{remark}[Finite Mostowski collapse]\label{rem:finitemostowski}
On rooted extensional presentations $E$ is an isomorphism (Lemma~\ref{lem:L2}); in general it
collapses. In Figure~\ref{fig:packaging} a five-node presentation is packaged to the three-node
canonical presentation of $2$. Nodes that cannot be reached from the point do not affect its
decoration and are discarded by $E$; on the part reachable from the point, $E$ is the finite case
of the Mostowski collapse. The infinitary statement is classical and is neither reproved nor used
here.
\end{remark}

\subsection{Closing from nothing yields exactly \texorpdfstring{$\HF$}{HF}}\label{ssec:hf}

By Theorem~\ref{thm:p5} we may work directly with the packaged objects, the elements of $\HF$. We
write the basic operations on $\HF$ as packaging clauses and show that closing the empty family
under them produces all of $\HF$.

\begin{definition}[The finite packaging operator]\label{def:F}
For $A\subseteq\HF$ let
\[
  \begin{aligned}
  F(A):={}&\{\varnothing\}\cup A
  \;\cup\;\underbrace{\{\{a,b\}:a,b\in A\}}_{\text{pairing (P5)}}
  \;\cup\;\underbrace{\{\textstyle\bigcup a:a\in A\}}_{\text{union (P5)}}\\
  &\cup\;\underbrace{\{S:S\subseteq a,\ a\in A\}}_{\text{gating (P2)}}
  \;\cup\;\underbrace{\{h[a]:a\in A,\ h\colon a\to A\}}_{\text{transport (P1)}},
  \end{aligned}
\]
where $h[a]=\{h(t):t\in a\}$. Because every $a\in\HF$ is finite, every subset $S\subseteq a$ and
every function $h\colon a\to A$ is admitted; no definability restriction is needed at this level.
Definability restrictions reappear, and matter, in the forcing section (\S\ref{sec:forcing}).
\end{definition}

\begin{lemma}[$\clF$ is a closure operator; \inhouse]\label{lem:closure}
$F$ is extensive and monotone, and $\HF$ is $F$-closed. Hence for each $A\subseteq\HF$ the least
$F$-closed superset $\clF(A)=\bigcap\{B: A\subseteq B\subseteq\HF,\ F(B)\subseteq B\}$ exists,
and $\clF$ is a closure operator: extensive, monotone and idempotent
\citep{DaveyPriestley2002,Tarski1955}.
\end{lemma}

\begin{proof}
Extensivity is built in. For monotonicity, if $A\subseteq B$ then every output of a clause applied
to $A$ is an output of the same clause applied to $B$; in particular a map $h\colon a\to A$ is
also a map into $B$. $\HF$ is closed clause by clause, since each clause outputs a finite set of
hereditarily finite sets. The intersection of $F$-closed families containing $A$ is again
$F$-closed, by monotonicity, and the closure properties follow.
\end{proof}

\noindent Although $\HF$ is countable, the lattice $\powset(\HF)$ in which this intersection is
taken is uncountable. Forming the intersection uses the ambient set-theoretic metatheory (with no
choice principle), not a finite computation. The same is true of every statement below about
$\HF$ as a whole: each object and each operation is finite, but the global statements are
statements of the metatheory.

\begin{lemma}[Finite collection; \inhouse]\label{lem:collection}
If $A$ is $F$-closed and $a_1,\dots,a_n\in A$, then $\{a_1,\dots,a_n\}\in A$.
\end{lemma}

\begin{proof}
Induction on $n$, using $x\cup\{y\}=\bigcup\{x,\{y\}\}$: the cases $n\le 1$ use the empty-set
and pairing clauses, and the inductive step pairs $\{a_1,\dots,a_n\}$ with $\{a_{n+1}\}$ and takes
the union.
\end{proof}

\noindent Finite collection is thus \emph{derived} from pairing and union, not assumed.

\begin{theorem}[$\HF$ is the saturation point of finite packaging; \inhouse]\label{thm:hfclosure}
$\clF(\varnothing)=\HF$.
\end{theorem}

\begin{proof}
$\clF(\varnothing)\subseteq\HF$ because $\HF$ is $F$-closed. For the reverse inclusion, let
$C=\clF(\varnothing)$ and induct on rank: $\varnothing\in C$, and if every member of $a\in\HF$ lies
in $C$, then $a=\{t_1,\dots,t_n\}$ is a finite collection of elements of $C$, so $a\in C$ by
Lemma~\ref{lem:collection}.
\end{proof}

\begin{corollary}[Finite power sets are reachable; \inhouse]\label{cor:finpow}
For $a\in\HF$, $\powset(a)\in\clF(\varnothing)$: each $S\subseteq a$ enters by the gate clause,
and $\powset(a)$, a finite collection of such sets, enters by Lemma~\ref{lem:collection}.
\end{corollary}

\noindent Corollary~\ref{cor:finpow} marks the boundary on which the paper turns. \emph{Finite}
power sets are free, a consequence of gating and finite collection. Collecting \emph{all}
subsets of an infinite set into one object is a different act, and it must be purchased
(\S\ref{sec:purchases}).

\begin{theorem}[$(\HF,\in)$ satisfies $\mathsf{ZFC}$ without Infinity; \inhouse]\label{thm:hfsat}
The structure $(\HF,\in)$ satisfies Extensionality, Empty set, Pairing, Union, Power set, every
instance of the Separation and Replacement schemas, Foundation and Choice, and it satisfies the
negation of Infinity.
\end{theorem}

\begin{proof}
$\HF$ is transitive, so for $a,b\in\HF$ the statements ``$a\in b$'' and ``$a\subseteq b$'' mean
the same inside $(\HF,\in)$ as outside, and every subset of an element of $\HF$ is again in
$\HF$. Extensionality follows from transitivity. Empty set, Pairing, Union and Power set are
supplied by the corresponding clauses of $F$ and by Corollary~\ref{cor:finpow}.
\emph{Separation:} for $a\in\HF$ and a formula $\varphi$ with parameters in $\HF$, the set
$\{x\in a:(\HF,\in)\models\varphi(x)\}$ is a subset of the finite set $a$, hence in $\HF$ by the
gate clause. Only truth in $(\HF,\in)$ is used; no absoluteness of unbounded formulas is needed.
\emph{Replacement:} if $\varphi$ defines in $(\HF,\in)$ a function on $a\in\HF$, its unique values
form a finite function $h\colon a\to\HF$, and $h[a]\in\HF$ by the transport clause.
\emph{Foundation:} a nonempty $a\in\HF$ is finite, so it has an element of minimal rank, which is
$\in$-minimal by Lemma~\ref{lem:L3}. \emph{Choice:} a finite family of nonempty sets in $\HF$ has
a choice function in $\HF$; picking the member with the least Ackermann code is one explicit
choice, and the resulting finite set of ordered pairs is in $\HF$. \emph{Infinity fails:} let
$a\in\HF$ contain $\varnothing$, and let $x\in a$ have maximal rank. Then $x\cup\{x\}$ has rank
$\rank(x)+1$, so it is not in $a$, and $a$ is not inductive.
\end{proof}

\begin{remark}[Arithmetic]\label{rem:ackermann}
$(\HF,\in)$ is the classical model $\Vfin$ of $\mathsf{ZF}$ with Infinity replaced by its
negation \citep{Kunen2011}. Ackermann's coding \citep{Ackermann1937} gives mutually
inverse interpretations between $(\HF,\in)$ and the standard model of arithmetic. At the level of
theories, bi-interpretability with Peano arithmetic requires finite set theory to include, besides
the negation of Infinity, the statement that every set has a transitive closure
\citep{KayeWong2007}; $\HF$ satisfies that statement, so the remark applies to the model
studied here.
\end{remark}

\subsection{Foundation is rank-audit exactness}\label{ssec:foundation}

The axiom of Foundation goes back to the well-foundedness analyses of \citet{Mirimanoff1917} and
\citet{VonNeumann1928}. On finite graphs it is equivalent to the existence of an audit.

\begin{theorem}[Finite rank-audit equivalence; \inhouse]\label{thm:tfae}
For a finite directed graph $G=(V,\to)$ the following are equivalent:
\begin{enumerate}\itemsep0pt
  \item $G$ is acyclic (it has no directed cycle, and in particular no self-loop);
  \item there is an audit $\rho\colon V\to\mathbb N$ that strictly decreases along edges:
  $u\to w$ implies $\rho(w)<\rho(u)$;
  \item every nonempty $S\subseteq V$ has a $\to$-minimal element.
\end{enumerate}
\end{theorem}

\begin{proof}
$(1)\Rightarrow(2)$: in an acyclic graph every path has at most $|V|-1$ edges, so the height
of Lemma~\ref{lem:L1} (the length of the longest path from a node) is defined and strictly
decreases along edges. $(2)\Rightarrow(1)$: a cycle $u_0\to\cdots\to u_{k-1}\to u_0$ would give
$\rho(u_0)<\rho(u_0)$; a self-loop is the case $k=1$. $(2)\Rightarrow(3)$: an element of $S$
with least $\rho$-value is $\to$-minimal. $(3)\Rightarrow(1)$: the node set of a cycle is a
nonempty set without a $\to$-minimal element.
\end{proof}

\noindent Natural-number ranks always suffice for finite acyclic graphs. Infinite well-founded
relations may require ordinal-valued ranks.

\begin{corollary}[\inhouse]\label{cor:foundation}
On a canonical presentation $C(a)$ the rank is a strictly decreasing audit (Lemma~\ref{lem:L3}),
so by Theorem~\ref{thm:tfae} every nonempty $b\subseteq\TC(\{a\})$ has an $\in$-minimal element.
Conversely, any finite presentation that admits a strictly decreasing audit is acyclic.
\end{corollary}

\begin{remark}[Foundation as a certificate]\label{rem:foundation}
Theorem~\ref{thm:tfae} reads the prohibition of membership cycles as a certificate that a
staging discipline is in force. Membership graphs with cycles---the self-loop $x=\{x\}$, the
two-cycle $x=\{y\}$, $y=\{x\}$---are perfectly coherent finite objects. They are excluded by
\emph{this} instrument because they carry no audit; anti-foundation \citep{Aczel1988} admits them.
This is countermodel \textbf{CM-2} (\S\ref{sec:atlas}).
\end{remark}

\subsection{The free-part ledger}

\begin{table}[tbp]\centering
\small
\begin{tabular}{@{}lll@{}}
\toprule
Axiom in $(\HF,\in)$ & Supplier & Role \\
\midrule
Extensionality & lens quotient (Theorem~\ref{thm:p5}) & packaging boundary \\
Empty set & seed clause $\{\varnothing\}$ & P5 \\
Pairing & pairing clause & P5 \\
Union & union clause & P5 \\
Separation & gate clause & P2 \\
Power set (of finite sets) & gates $+$ finite collection (Corollary~\ref{cor:finpow}) & P2 $+$ P5 \\
Replacement (on finite domains) & transport clause & P1 \\
Choice (for finite families) & least-code selection & --- \\
Foundation & rank audit (Theorem~\ref{thm:tfae}) & P6 \\
$\neg\,$Infinity & every element is finite (Theorem~\ref{thm:hfsat}) & boundary \\
\bottomrule
\end{tabular}
\caption{The free part. Each axiom holds in $(\HF,\in)$ and is witnessed by a named clause of the
finite packaging operator $F$, by finite selection, or by the rank audit; the closure of the
empty family under these clauses is exactly $\HF$ (Theorem~\ref{thm:hfclosure}). All objects and
operations are finite; global statements about $\HF$ are made in the ambient metatheory.}
\label{tab:freepart}
\end{table}

Every entry of Table~\ref{tab:freepart} is proved here, and several finite cores behind the
entries are verified in Lean; the schema-level satisfaction statements remain on paper
(\S\ref{sec:mech}). The next two sections show that the boundary of this column is sharp:
Infinity cannot be reached from inside it (\S\ref{sec:infinity}), and forgetting the staging
loses nothing (\S\ref{sec:shadow}).
